% Intro

\vfill{}
\noindent\textbf{To appear in the Proceedings of the ACM SIGOPS 31st Symposium on Operating Systems Principles (SOSP ’25), October 13–16, 2025, Seoul, Republic of Korea.}

\section{Introduction}

Deep Learning~(DL) has had enormous success in domains such as robotics~\cite{hwangbo2019learning}, games~\cite{silver2017mastering}, natural language~\cite{brown2020language}, protein folding~\cite{jumper2021highly}, and image generation~\cite{rombach2022high}. DL computation is defined using operators~(\eg matrix multiplication, addition, and indexing) with input and output \emph{tensors}. Tensors are multi-dimensional arrays laid out spatially in memory, whose \emph{shape}---a tuple of integers---specifies the size of each spatial dimension. DL workloads are resource-intensive, and DL systems, such as Torch~\cite{paszke2019pytorch} or JAX~\cite{frostig2018compiling}, offer high-level tensor programming models for GPU accelerators~\cite{raina2009large}. 

A common feature of many DL algorithms~\cite{vaswani2017attention, sutton2018reinforcement, rumelhart1986learning, rombach2022high, gu2021efficiently} is their mathematical definition through \emph{recurrence equations}~\cite{karp1967organization}. Recurrence equations express \emph{temporal} relationships between values at different timesteps (past and future) declaratively, \ie without specifying an order of execution. For example, attention~\cite{vaswani2017attention}---a key operation in large language models~(LLMs)---depends, at a given timestep, on outputs from all prior timesteps; reinforcement learning~(RL)~\cite{sutton1988learning, williams1992simple} algorithms produce actions at each timestep and use a sum of future rewards to evaluate them.


Such \emph{dynamic} (\ie time-varying) dependencies between operators and tensors of different timesteps  introduce challenges when existing DL systems try to support them:
%
\begin{myitemize}
  \item \textbf{Expressiveness}---dynamic dependencies can create tensors with \emph{dynamic shapes} whose support is limited in terms of current APIs~\cite{ansel2024pytorch,lai2025relax}; 
  \item \textbf{Optimization}---existing graph representations~\cite{abadi2016tensorflow, frostig2018compiling} cannot express dynamic dependencies, missing opportunities for algebraic simplification or vectorization;  
  \item \textbf{Scheduling}---future dependencies complicate scheduling, as dependent computations must be delayed~\cite{karp1967organization, moritz2018ray};
  \item \textbf{Memory management}---dependencies on past timesteps result in complex decisions on which tensors to retain, deallocate or evict with limited GPU memory~\cite{rhu2016vdnn}.
\end{myitemize}

\begin{figure*}[t]
  \centering
  \includegraphics[width=1.0\textwidth]{cr_figures/overview}
  \caption{Overview of how \sys{} expresses, optimizes, schedules and manages the memory of dynamic DL algorithms}
  \label{fig:overview}
\end{figure*}

\noindent
\emph{Eager} DL systems, such as Torch~\cite{paszke2019pytorch} or TF-Eager~\cite{agrawal2019tensorflow}, support dynamism through the host language~(Python) by dispatching GPU kernels on-demand~\cite{chetlur2014cudnn}. This imperative approach, however, prevents them from performing compiler-based optimizations~\cite{ansel2024pytorch} or scheduling; in contrast, \emph{graph-based systems}, such as TensorFlow~\cite{abadi2016tensorflow} and JAX~\cite{frostig2018compiling}, enable ahead-of-time compilation by first constructing a \emph{dataflow graph} in which each node represents a DL operator and edges are tensors that flow between them. 
Optimizations such as memory planning~\cite{maas2022telamalloc,liu2024uncovering}, operator fusion~\cite{chen2018tvm}, vectorization~\cite{ragan2013halide}, and code-generation~\cite{Lattner_2021, vasilache2018tensor}, however, require statically known fixed-size tensor shapes.

To compile dynamic DL algorithms, users must \emph{pad} dynamic tensor shapes to a static size, and use \emph{masks} to encode dynamic dependencies, which leads to large memory and execution overheads~\cite{zheng2023bladedisc}. Alternatively, users can break-up computation into static sub-programs, but this prevents whole-program optimizations leading to worse performance~\cite{hessel2021podracer}.

Our key insight is that tensors in dynamic DL algorithms have implicit \emph{temporal dimensions}.
By making them explicit, dynamic dependencies become regular tensor-indexing operations across temporal dimensions. Furthermore, we show that dynamism in dependencies and shapes has structure, and can be concisely represented using \emph{symbolic expressions} of timestep symbols, enabling optimization.

We describe \sys{}, a new DL system that expresses, optimizes, and schedules dynamic DL programs, as well as manages their memory. \sys{} makes the following novel technical contributions (shown in \F\ref{fig:overview}):

\mypar{(1)~Declarative programming with recurrent tensors~(\S\ref{sec:rts})} To capture dynamic dependencies, we introduce \emph{recurrent tensors}~(RTs), a declarative programming model inspired by recurrence equations~\cite{karp1967organization}, often used in DL algorithms~\cite{schulman2015high, vaswani2017attention, sutton1988learning, rumelhart1986learning}. RTs can have multiple temporal dimensions (\eg timestep and iteration), which define a temporal \emph{domain}. RTs can be indexed using symbolic expressions to declaratively express dynamic dependencies. RTs have well-defined domain broadcast~\cite{van2011numpy} and backpropagation~\cite{baydin2018automatic} semantics. 

\mypar{(2)~Optimization using symbolic dependence graphs~(\S\ref{sec:pdg_and_transformations})} \sys translates RTs into a \emph{symbolic dependence graph~(SDG)} in which each vertex operator carries the temporal domain of the RT that created it, \ie the set of timesteps at which the operator produces output tensors; edges describe how operators at different timesteps \emph{depend on} each other's outputs using the symbolic expressions extracted from RTs. 

SDGs enable \sys{} to extend classic compiler optimizations (\eg algebraic simplifications, redundant code elimination) to dynamic computations by manipulating symbolic expressions. In addition, SDGs simplify advanced transformations, such as \emph{vectorization} (batching many small operations) and \emph{tiling} (executing large operations in smaller tiles), by viewing them as moving tensor dimensions from temporal to spatial and vice versa. For efficient code-generation, \sys{} tiles dynamic dependencies into static-size tiles and \emph{fuses} islands of static operations into single operations.


\mypar{(3)~Scheduling using a polyhedral model~(\S\ref{sec:polyhedral_model})} The declarative RT program does not prescribe an execution order, and, due to their dynamic dependencies, SDGs cannot be scheduled using  topological sorting. Instead, \sys{} uses the polyhedral model~\cite{verdoolaege2010isl, bondhugula2008pluto} for program-wide scheduling: \sys{} translates the SDG into a set of constraints and solves an integer linear program to obtain a schedule that assigns a physical execution time to each timestep. Based on the schedule, \sys{} generates an imperative program for the user computation, represented as an abstract syntax tree~(AST).


\mypar{(4)~Automated memory management~(\S\ref{sec:mem_sched})} \sys{} introduces a simple method for automatic memory management, including buffer deallocations, donations, and GPU/CPU buffer swapping. It augments the SDG with explicit memory management operations, which are restricted by dependency edges. Therefore, memory management is also handled when solving for an execution schedule.


\tinyskip

\noindent
Our prototype implementation of \sys{}\footnote{Available  at \url{https://github.com/LSDS/Tempo}} supports single-GPU execution on multiple execution engines (Torch and JAX). To execute SDGs, \sys{} picks physical tensor memory representations based on dependence expressions and applies existing code-generators~\cite{openxla2024} to fused operators.
It then follows the schedule AST, evaluating symbolic dependence expressions, to resolve dynamic dependencies at runtime.

Our evaluation shows that \sys{} effectively compiles dynamic DL programs, while automating algorithm-specific scheduling and memory management. For LLM decoding using Llama-3.2-3B~\cite{dubey2024llama}, we compare \sys{} against Torch and JAX and show that it is up to 7$\times$ faster and scales to 4$\times$ longer sequences due to its awareness of dynamic dependencies. For three typical RL algorithms~\cite{schulman2017proximal, williams1992simple, sutton1988learning} with different dynamic dependencies, \sys achieves up to 54$\times$ faster end-to-end training times than five baseline RL systems by eliminating redundant computation. For both workloads, \sys{} has an up to 16$\times$ lower peak GPU memory usage due to its tiling optimization and efficient memory management.

% End

